\documentclass[11pt,a4paper]{article}
\usepackage[utf8]{inputenc}
\usepackage[T1]{fontenc}
\usepackage[english]{babel}
\usepackage{amsmath, amsthm, amssymb, amsfonts}
\usepackage{enumerate}
\usepackage{graphicx}
\usepackage{fancyhdr}
\usepackage{booktabs}
\usepackage{float}
\usepackage[colorlinks=true,breaklinks=false]{hyperref}
\usepackage{listings}
\usepackage[norelsize,ruled,vlined]{algorithm2e}
\usepackage{csquotes}
\usepackage{url}
\usepackage{ifthen}
\usepackage{keyval}
\usepackage{etoolbox}
\usepackage{etex}
\usepackage{datetime}
\usepackage{eurosym}
\usepackage{tikz}
\usepackage{wrapfig}
\usepackage{paralist}
\usepackage[backend=bibtex]{biblatex}
\usepackage{xcolor}
\usepackage{caption}
\DeclareCaptionFont{white}{\color{white}}
\DeclareCaptionFormat{listing}{\colorbox[cmyk]{0.43, 0.35, 0.35,0.01}{\parbox{\textwidth}{\hspace{15pt}#1#2#3} } }
\captionsetup[lstlisting]{format=listing, labelfont=white, textfont=white, singlelinecheck=false, margin=0pt, font={bf,footnotesize} }

\addbibresource{project.bib}

% Single spacing after dots and colons
\frenchspacing
% \hyphenation{directembedding}
% \hyphenation{DSLVirtualization}
% \hyphenation{ReificationTransformation}

%%%%%%%%%%%%%%%%%%%%%%%%%%%%%%%%%%%%%%%%%%%%%%%%%%%%%%%%%%%%%
%                        Setup
%%%%%%%%%%%%%%%%%%%%%%%%%%%%%%%%%%%%%%%%%%%%%%%%%%%%%%%%%%%%%

\begin{document}

\definecolor{dkgreen}{rgb}{0,0.6,0}
\definecolor{gray}{rgb}{0.5,0.5,0.5}
\definecolor{mauve}{rgb}{0.58,0,0.82}

\lstset{language=Scala,
  aboveskip=3mm,
  basicstyle={\small\ttfamily},
  belowskip=3mm,
  breakatwhitespace=true,
  breaklines=true,
  columns=flexible,
  commentstyle=\color{dkgreen},
  keywordstyle=\color{blue},
  numbers=left,
  numberstyle=\tiny\color{gray},
  showstringspaces=false,
  stepnumber=1,
  stringstyle=\color{mauve},
  tabsize=2,
}

\usetikzlibrary{calc,trees,positioning,arrows,chains,shapes.geometric,%
    decorations.pathreplacing,decorations.pathmorphing,shapes,%
    matrix,shapes.symbols}
\tikzset{>=stealth',
  punktchain/.style={rectangle,
    rounded corners,
    % fill=black!10,
    draw=black, very thick,
    text width=12.5em,
    minimum height=2.5em,
    text centered,
    on chain},
  every join/.style={->, shorten >=2pt},
}


% Example of title page for the projects carried out within the lasec

% Simply include it in your mastex tex file:
%        % Example of title page for the projects carried out within the lasec

% Simply include it in your mastex tex file:
%        % Example of title page for the projects carried out within the lasec

% Simply include it in your mastex tex file:
%        \input{cover}


% Updated March 2006 (SP)


\newcommand{\logoepfl}[0]{
  \begin{center}
    \includegraphics[width=4cm]{logo_epfl_coul.eps}
  \end{center}
  \vspace{0.3cm}
  \hrule
}
\newcommand{\logolasec}[0]{
  \vspace{1cm}
  \hrule
  \begin{center}
    \includegraphics[width=4.5cm]{logo_lasec_coul.eps}
  \end{center}
}
\newcommand{\project}[1]{
  \begin{center}
    \large{#1}
  \end{center}
  \vspace{1cm}
}
\newcommand{\department}[1]{
  \begin{center}
    \large{#1}
  \end{center}
}
\newcommand{\supervisor}[3]{
  \begin{center}
    \begin{normalsize}{
        \bf #1}\\#2\\#3
    \end{normalsize}
  \end{center}
}
\renewcommand{\author}[1]{
  \begin{center}
    \Large{#1}
  \end{center}
  \vspace{0.5cm}
}
\renewcommand{\title}[1]{
  \vspace{3cm}
  \begin{center}
    \huge{#1}
  \end{center}
  \vspace{1.7cm}
}
\renewcommand{\date}[2]{
  \begin{center}
    \normalsize{#1 #2}
  \end{center}
  \vspace{0.5cm}
}


\thispagestyle{empty}


% begin title page
  \logoepfl

  \title{scalafmt: opinionated \\ code formatter for Scala}

  \author{Ólafur Páll Geirsson}
  \department{School of Computer and Communication Sciences}
  \project{A thesis submitted for the degree of Master of Computer Science at École polytechnique fédérale de Lausanne}

  \date{June}{2016}

  \begin{center}
    \begin{tabular}{cc}
      \begin{tabular}{p{4.0cm}}
        \supervisor{Responsible}{Prof. Martin Odersky}{EPFL / LAMP}
      \end{tabular}&
      \begin{tabular}{p{4.0cm}}
        \supervisor{Supervisor}{Eugene Burmako}{EPFL / LAMP}
      \end{tabular}
    \end{tabular}
  \end{center}

% end title page




% Updated March 2006 (SP)


\newcommand{\logoepfl}[0]{
  \begin{center}
    \includegraphics[width=4cm]{logo_epfl_coul.eps}
  \end{center}
  \vspace{0.3cm}
  \hrule
}
\newcommand{\logolasec}[0]{
  \vspace{1cm}
  \hrule
  \begin{center}
    \includegraphics[width=4.5cm]{logo_lasec_coul.eps}
  \end{center}
}
\newcommand{\project}[1]{
  \begin{center}
    \large{#1}
  \end{center}
  \vspace{1cm}
}
\newcommand{\department}[1]{
  \begin{center}
    \large{#1}
  \end{center}
}
\newcommand{\supervisor}[3]{
  \begin{center}
    \begin{normalsize}{
        \bf #1}\\#2\\#3
    \end{normalsize}
  \end{center}
}
\renewcommand{\author}[1]{
  \begin{center}
    \Large{#1}
  \end{center}
  \vspace{0.5cm}
}
\renewcommand{\title}[1]{
  \vspace{3cm}
  \begin{center}
    \huge{#1}
  \end{center}
  \vspace{1.7cm}
}
\renewcommand{\date}[2]{
  \begin{center}
    \normalsize{#1 #2}
  \end{center}
  \vspace{0.5cm}
}


\thispagestyle{empty}


% begin title page
  \logoepfl

  \title{scalafmt: opinionated \\ code formatter for Scala}

  \author{Ólafur Páll Geirsson}
  \department{School of Computer and Communication Sciences}
  \project{A thesis submitted for the degree of Master of Computer Science at École polytechnique fédérale de Lausanne}

  \date{June}{2016}

  \begin{center}
    \begin{tabular}{cc}
      \begin{tabular}{p{4.0cm}}
        \supervisor{Responsible}{Prof. Martin Odersky}{EPFL / LAMP}
      \end{tabular}&
      \begin{tabular}{p{4.0cm}}
        \supervisor{Supervisor}{Eugene Burmako}{EPFL / LAMP}
      \end{tabular}
    \end{tabular}
  \end{center}

% end title page




% Updated March 2006 (SP)


\newcommand{\logoepfl}[0]{
  \begin{center}
    \includegraphics[width=4cm]{logo_epfl_coul.eps}
  \end{center}
  \vspace{0.3cm}
  \hrule
}
\newcommand{\logolasec}[0]{
  \vspace{1cm}
  \hrule
  \begin{center}
    \includegraphics[width=4.5cm]{logo_lasec_coul.eps}
  \end{center}
}
\newcommand{\project}[1]{
  \begin{center}
    \large{#1}
  \end{center}
  \vspace{1cm}
}
\newcommand{\department}[1]{
  \begin{center}
    \large{#1}
  \end{center}
}
\newcommand{\supervisor}[3]{
  \begin{center}
    \begin{normalsize}{
        \bf #1}\\#2\\#3
    \end{normalsize}
  \end{center}
}
\renewcommand{\author}[1]{
  \begin{center}
    \Large{#1}
  \end{center}
  \vspace{0.5cm}
}
\renewcommand{\title}[1]{
  \vspace{3cm}
  \begin{center}
    \huge{#1}
  \end{center}
  \vspace{1.7cm}
}
\renewcommand{\date}[2]{
  \begin{center}
    \normalsize{#1 #2}
  \end{center}
  \vspace{0.5cm}
}


\thispagestyle{empty}


% begin title page
  \logoepfl

  \title{scalafmt: opinionated \\ code formatter for Scala}

  \author{Ólafur Páll Geirsson}
  \department{School of Computer and Communication Sciences}
  \project{A thesis submitted for the degree of Master of Computer Science at École polytechnique fédérale de Lausanne}

  \date{June}{2016}

  \begin{center}
    \begin{tabular}{cc}
      \begin{tabular}{p{4.0cm}}
        \supervisor{Responsible}{Prof. Martin Odersky}{EPFL / LAMP}
      \end{tabular}&
      \begin{tabular}{p{4.0cm}}
        \supervisor{Supervisor}{Eugene Burmako}{EPFL / LAMP}
      \end{tabular}
    \end{tabular}
  \end{center}

% end title page


\begin{abstract}
  Automatic code formatters bring a lot of benefits to software development.
  When done right, code formatters relieve the developer's attention from
  syntactic trivia and enforce a consistent coding style between developers.
  However, developing a good code formatter is still somewhat of a black art and yet,
  little research has been made towards the algorithms that power such code formatters.

  This thesis addresses the problem of developing a feature-rich code formatter for a custom programming language.
  % A critical feature of such a formatter is enforcing a column-width limit, which turns out to be a hard problem.
  Our contributions are twofold.
  First, we present a language agnostic framework that consists of core data structures, algorithms and tooling 
  that allow rapid development and testing of such a code formatter.
  Secondly, we use our model in a case study and provide a concrete implementation of a code formatter, \texttt{scalafmt}.
  We show that the framework makes


\end{abstract}

\begin{abstract}
  Automatic code formatters bring many benefits to software development, yet they can be tricky to get right.
  This thesis addresses the problem of developing a code formatter for the
  Scala programming language that captures the language's most popular idioms
  and coding styles.
  Our work has been limited to formatting Scala code.
  Still, we have developed data structures, algorithms and tools that we
  believe can applied to develop code formatters for a variety of other
  programming languages.
\end{abstract}
\tableofcontents
% 
\section{Introduction} % (fold)
\label{sec:Introduction}
The main contributions presented in this thesis are the following:
\begin{itemize}
    \item A DSL for declaring feature-rich code formatting rules.
    \item An case study where the DSL is used to format Scala programs.
\end{itemize}
Throughout the paper we assume familiarity with the basics of the Scala
Programming Language~\autocite{odersky_scala_2004}.

\section{Introduction} % (fold)
\label{sec:Introduction}
Software developers use code formatters to make their code look pretty.
A code formatter should not change the behaviour of code.
Typically, a code formatter only modifies whitespace tokens, or syntactic trivia.
Syntactic trivia has nothing to do with the behavior of code.
Code formatters promise to relieve software developers from manually manipulating
syntactic trivia.
Instead, developers can focus on writing correct, performant and maintainable code.

Code formatters brings many other benefits to software development.
They can help enforce a consistent coding style in a codebase that's touched by multiple developers.
A consistent code style aids readability of code.
Code formatting also enables automated refactoring tools, see~\autocite{wright_large-scale_2013}.
Code formatting also lowers the barrier to entry for novice programmers, who are not aware of the coding style conventions.

Code formatters have existed for a long time.
In more recent years we've seen many new opimization based code formatters, such as clang-format and dartfmt.
Optimization based code formatters define a fit function for a given layout.
A good layout has low cost while a bad layout has a high cost.
The formatter runs an algorithm to choose the layout with the lowest code.
This idea is not new, for example TeX uses a fit function to optimally break lines
in text documents.

Much work has been put into developing language agnostic code formatters or pretty printers.
Researchers of such work correctly highlights that developing a custom code formatter for each language requires a lot of effort.
It may be possible that we find a perfect way to develop language independent code formatters, see~\autocite{mps_article}.
Still, today no such framework exists and people who write Scala code each day are still choosing to manually format their code.
The goal of this thesis is not to address language agnostic code formatting.

\subsection{Problem statement}
This thesis addresses the problem of developing a code formatter, \texttt{scalafmt}, for the Scala programming language.
This formatter should:
\begin{itemize}
  \item capture Scala's most common idioms and popular coding styles, most importanly a line length limit.
  \item be opinionated, it should be able to produce nice looking code even from computer generated input.
  \item be fast, <200ms for 99 percentile of source files, <2s for rest.
\end{itemize}

\subsection{Contributions}
The main contributions presented in this thesis are the following:
\begin{itemize}
    \item We have developed a framework of data structures, algorithms and tooling that help
      developing code formatters. This work is presented in section X.
    \item We have applied this framework to develop \texttt{scalafmt}, a code
      formatter for the Scala programming language.
\end{itemize}

\section{Background}
\subsection{Scala the programming language}

Throughout the paper we assume familiarity with the basics of the Scala
Programming Language~\autocite{odersky_scala_2004}.
\subsubsection{Higher order functions}
\subsubsection{Pattern matching}
\subsubsection{Metaprogramming with scala.meta}
\begin{itemize}
  \item Tree nodes have parent links.
  \item Token classes are types.
\end{itemize}
\subsection{Code formatting}
\subsubsection{Scalariform}
\section{Framework}
\subsection{Data structures}
\subsubsection{FormatToken}
\subsubsection{Split combinator}
\subsubsection{Policy combinator}
\subsubsection{OptimalToken combinator}
\subsubsection{State}
\subsubsection{Indent}
\subsection{Algorithms}
Pipeline of three stages: pre-processing, line wrapping and post-processing.
\subsubsection{Pre-processing}
 \begin{itemize}
   \item AST vs. token stream
 \end{itemize}
rustfmt considers that all eventually converge to a hybrid of the two.
Indeed, clang-format has an elaborate parser.
\subsubsection{Layout}
\begin{itemize}
  \item Router
  \item Search.
\end{itemize}
\subsubsection{Post-processing}
\begin{itemize}
  \item vertical alignment
  \item comment formatting
  \item stripMargin alignment
\end{itemize}
\subsection{Tooling}
\subsubsection{Heatmaps}
\subsubsection{Unit tests}
\subsubsection{Property based testing}
\begin{enumerate}
  \item AST integrity
  \item Idempotent
\end{enumerate}
\section{scalafmt}
\subsection{Custom optimizations}
\begin{enumerate}
  \item dequeueOnNewStatements
  \item recurseOnBlocks
  \item escapeInPathologicalCases
  \item pruneSlowStates
\end{enumerate}
\section{Evaluation}
\subsection{Performance}
\subsection{Output}
\begin{itemize}
  \item gofmt
  \item rustfmt
  \item dartfmt
  \item clang-format
  \item scalariform
\end{itemize}
\section{Related work}
\subsection{Combinator based}
\begin{enumerate}
  \item Houghes 1995
  \item Wadler 1999
\end{enumerate}
\subsection{Optimization-oriented}
\begin{enumerate}
  \item clang-format Dijkstra's 2010
  \item dartfmt Best-first search 2014
  \item rfmt 2015
\end{enumerate}
\begin{enumerate}
  \item Optimal line breaking
  \item Oppen
\end{enumerate}
\section{Discussion}
\subsection{Future work}
\subsection{Conclusion}
\printbibliography{}

\end{document}

